\documentclass[10pt, aspectratio=169]{beamer}
\input{../../_en_preambulo_beamer}
% Comandos y macros estadísticas compartidas para las presentaciones Beamer en inglés (Metropolis)

% Conjuntos numéricos básicos
\newcommand{\R}{\mathbb{R}}
\newcommand{\N}{\mathbb{N}}
\newcommand{\Q}{\mathbb{Q}}
\newcommand{\Z}{\mathbb{Z}}
\newcommand{\set}[1]{\left\{ #1 \right\}}
\newcommand{\abs}[1]{\left|#1\right|}
\newcommand{\norm}[1]{\left\| #1 \right\|}

% Comandos estadísticos y probabilísticos del curso
\newcommand{\Var}{\operatorname{Var}}
\newcommand{\cov}{\operatorname{Cov}}
\newcommand{\corr}{\rho}
\newcommand{\comb}[2]{\begin{pmatrix} #1 \\ #2 \end{pmatrix}}
\newcommand{\s}{\sigma}
\newcommand{\particion}{\mathfrak{P}}
\newcommand{\card}[1]{\#\left( #1 \right)}
\newcommand{\seq}[3]{\set{#1}_{#2}^{#3}}
\newcommand{\E}{\mathbb{E}}
\newcommand{\Prob}{\mathbb{P}}

% Entornos pedagógicos y cajas en inglés académico natural (soportando tanto nombres en inglés como en español)
\newenvironment{discussion}{%
  \begin{alertblock}{Discussion Question}%
}{%
  \end{alertblock}%
}
\newenvironment{discusion}{%
  \begin{alertblock}{Discussion Question}%
}{%
  \end{alertblock}%
}

\newenvironment{reflection}{%
  \begin{alertblock}{Classroom Reflection}%
}{%
  \end{alertblock}%
}
\newenvironment{reflexion}{%
  \begin{alertblock}{Classroom Reflection}%
}{%
  \end{alertblock}%
}

\newenvironment{paradox}{%
  \begin{alertblock}{Exploratory Paradox}%
}{%
  \end{alertblock}%
}
\newenvironment{paradoja}{%
  \begin{alertblock}{Exploratory Paradox}%
}{%
  \end{alertblock}%
}

\newenvironment{motivation}{%
  \begin{exampleblock}{Motivation: Data Science in Practice}%
}{%
  \end{exampleblock}%
}
\newenvironment{motivacion}{%
  \begin{exampleblock}{Motivation: Data Science in Practice}%
}{%
  \end{exampleblock}%
}

\newenvironment{quickchallenge}{%
  \begin{block}{Quick Team Challenge}%
}{%
  \end{block}%
}
\newenvironment{retoclase}{%
  \begin{block}{Quick Team Challenge}%
}{%
  \end{block}%
}


\title[Correlation]{Section 09.01: Correlation as the Premise of Regression}
\subtitle{Definition, Properties, and the Correlation-Causation Warning}
\author[J. Castillo Colmenares]{Juliho Castillo Colmenares}
\institute{Tecnológico de Monterrey}
\date{\vspace{-1.5cm}}

\begin{document}

% Slide 1: Title
\begin{frame}[plain]
  \titlepage
\end{frame}

% Slide 2: Roadmap
\begin{frame}{Roadmap --- Chapter 09: Linear and Multiple Regressions}
  \scriptsize
  Where are we in the modular development of this chapter?
  \begin{itemize}\itemsep=0.02em
    \item \textbf{09.01} Correlation as the Premise of Regression \emph{(Today)}
    \item \textbf{09.02} Introduction to Linear Regression
    \item \textbf{09.03} The Mathematics of Regression: OLS Derivation and $R^2$
    \item \textbf{09.04} Linear Regression on Simulated Data
    \item \textbf{09.05} Optimal Coefficients, $t$/$F$ Tests, and RSE
    \item \textbf{09.06} Implementation in Python with \texttt{statsmodels}
    \item \textbf{09.07} Multiple Regression, the Normal Equation, and Ridge/Lasso
    \item \textbf{09.08} Model Validation and $k$-fold Cross-Validation
    \item \textbf{09.09} Residual Diagnostics and Multicollinearity (VIF)
    \item \textbf{09.10} Regression with \texttt{scikit-learn}
    \item \textbf{09.11} Categorical Variables and Dummy Variables
    \item \textbf{09.12} Nonlinear Transformations and Polynomial Regression
  \end{itemize}
  \pause
  \vspace{-0.05cm}
  \begin{block}{Session Objective}
    Establish correlation as the conceptual premise of all predictive modeling: formally define the Pearson coefficient, study its properties and range, and rigorously warn that correlation does \textbf{not} imply causation --- the distinction that opens the door to linear regression in Section 09.02.
  \end{block}
\end{frame}

% Slide 3: Motivation
\begin{frame}{Why Start with Correlation?}
  \footnotesize
  \begin{motivacion}
    Correlation is the premise of predictive modeling: it is the factor on which we can anticipate the behavior of one variable from another. If $X$ and $Y$ are strongly correlated, a change in one suggests an associated change in the other.
  \end{motivacion}

  \vspace{0.15cm}
  \pause
  \begin{itemize}\itemsep=0.1cm
    \item A good correlation guarantees an approximate functional mathematical relationship $Y=f(X)$. \pause
    \item That relationship can take any form: linear ($y=ax+b$), exponential ($y=be^{ax}$), or other --- linear regression focuses on the linear case.
  \end{itemize}
\end{frame}

% Slide 4: Formal definition
\begin{frame}{Formal Definition of the Correlation Coefficient}
  \footnotesize
  \begin{block}{Pearson's Correlation Coefficient}
    For variables $X,Y$ with sample means $\bar x,\bar y$:
    $$r_{XY}=\frac{\sum_i(x_i-\bar x)(y_i-\bar y)}{\sqrt{\sum_i(x_i-\bar x)^2\sum_i(y_i-\bar y)^2}}=\frac{S_{XY}}{\sqrt{S_{XX}S_{YY}}}$$
  \end{block}
  \pause

  \vspace{0.1cm}
  \begin{alertblock}{Interpretation}
    $r_{XY}$ normalizes the covariance $S_{XY}$ by the individual dispersions of $X$ and $Y$, producing a dimensionless measure of linear association comparable across any pair of variables.
  \end{alertblock}
\end{frame}

% Slide 5: Properties
\begin{frame}{Properties of the Correlation Coefficient}
  \footnotesize
  \begin{itemize}\itemsep=0.12cm
    \item \textbf{Bounded range:} $-1\leq r_{XY}\leq1$, a direct consequence of the Cauchy-Schwarz inequality. \pause
    \item \textbf{Sign:} $r_{XY}>0$ indicates a direct relationship (both variables rise or fall together); $r_{XY}<0$ indicates an inverse relationship. \pause
    \item \textbf{Magnitude:} the larger $|r_{XY}|$, the stronger and more predictable the linear relationship; $|r_{XY}|=1$ indicates an exact linear relationship.
  \end{itemize}
\end{frame}

% Slide 6: Warning
\begin{frame}{Warning! Correlation Does Not Imply Causation}
  \footnotesize
  \begin{alertblock}{Critical Observation}
    A strong correlation suggests a relationship useful for prediction, but \textbf{does not imply that relationship is the only factor} explaining the observed behavior, nor that one variable \emph{causes} the other.
  \end{alertblock}
  \pause

  \vspace{0.1cm}
  \begin{block}{Classic Example: the Confounding Variable}
    There is a very strong positive correlation between the number of words a child knows and their shoe size. Neither causes the other: both are consequences of a third hidden variable, \textbf{age}.
  \end{block}
\end{frame}

% Slide 7: Advertising case study
\begin{frame}{Case Study: Advertising and Sales}
  \footnotesize
  Classic dataset: advertising spend by medium (\texttt{TV}, \texttt{Radio}, \texttt{Newspaper}) vs. \texttt{Sales} of a product.
  \pause

  \vspace{0.1cm}
  \begin{block}{Observed Correlation Coefficients}
    $$r(\texttt{TV},\texttt{Sales})\approx0.782, \qquad r(\texttt{Radio},\texttt{Sales})\approx0.576$$
  \end{block}
  \pause

  \vspace{0.1cm}
  \begin{alertblock}{Preliminary Reading}
    \texttt{TV} shows the strongest linear association with \texttt{Sales}: a natural candidate for the first predictor of a regression model (Sections 09.02 onward).
  \end{alertblock}
\end{frame}

% Slide 8: Python Lab (Block 1)
\begin{frame}[fragile]{Python Lab: Pearson's Coefficient from Its Definition}
  \scriptsize
  Direct computation of $r_{XY}$ from sums of deviations, verified against \texttt{np.corrcoef} (\texttt{09.01\_correlation.py}):
  {\lstset{basicstyle=\fontsize{3.6pt}{4.3pt}\selectfont\ttfamily,aboveskip=0.05em,belowskip=0.05em}
  \lstinputlisting[firstline=17, lastline=27]{../../code/09_regresiones/09.01_correlation.py}}
\end{frame}

% Slide 9: Python Lab (Block 2)
\begin{frame}[fragile]{Python Lab: Range and Strength of Correlation}
  \scriptsize
  Six scenarios illustrating the $[-1,1]$ range and sign/magnitude interpretation:
  {\lstset{basicstyle=\fontsize{3.4pt}{4.1pt}\selectfont\ttfamily,aboveskip=0.05em,belowskip=0.05em}
  \lstinputlisting[firstline=30, lastline=45]{../../code/09_regresiones/09.01_correlation.py}}
\end{frame}

% Slide 10: Python Lab (Block 3)
\begin{frame}[fragile]{Python Lab: Correlation Is Not Causation}
  \scriptsize
  Spurious correlation between shoe size and vocabulary, resolved via partial correlation controlling for age:
  {\lstset{basicstyle=\fontsize{3.2pt}{3.9pt}\selectfont\ttfamily,aboveskip=0.05em,belowskip=0.05em}
  \lstinputlisting[firstline=48, lastline=64]{../../code/09_regresiones/09.01_correlation.py}}
\end{frame}

% Slide 11: Terminal Output
\begin{frame}[fragile]{Python Lab: Terminal Output}
  \scriptsize
  Standard output produced when running the lab script:
  \vspace{0.02cm}
  \begin{block}{Standard Console Output}
    \fontsize{3.8pt}{4.6pt}\selectfont
    \begin{verbatim}
=== Block 1: Pearson from Definition ===
r (from definition) = 0.8959 (matches np.corrcoef)

=== Block 2: Range and Strength ===
perfect positive: r=+1.0000   strong positive: r=+0.9638
weak positive:     r=+0.3710   no relationship: r=-0.1333
strong negative:   r=-0.9690   perfect negative: r=-1.0000

=== Block 3: Correlation is Not Causation ===
Raw correlation(shoe size, vocabulary) = 0.9517 (spurious)
Partial correlation(... | age)         = 0.0486 (collapses)
    \end{verbatim}
  \end{block}
\end{frame}

% Slide 12: In-class discussion (Statement)
\begin{frame}{In-Class Discussion --- The Danger of a Single Coefficient}
  \footnotesize
  \begin{block}{Question}
    In the advertising case, $r(\texttt{Radio},\texttt{Sales})\approx0.576$ is a moderate, not weak, correlation. Why isn't looking at $r$ alone enough to decide whether \texttt{Radio} should be included in a multiple regression model together with \texttt{TV}?
  \end{block}

  \vspace{0.15cm}
  \pause
  \begin{alertblock}{Hint}
    Think about whether \texttt{Radio} and \texttt{TV} might themselves be correlated with each other.
  \end{alertblock}
\end{frame}

% Slide 13: In-class discussion (Answer)
\begin{frame}{In-Class Discussion --- Answer}
  \scriptsize
  The bivariate correlation $r(\texttt{Radio},\texttt{Sales})$ completely ignores the possible relationship between \texttt{Radio} and \texttt{TV}: \pause

  \vspace{0.1cm}
  \begin{alertblock}{Interpretation}
    If \texttt{Radio} and \texttt{TV} share variability (for example, campaigns combining both media simultaneously), part of what $r(\texttt{Radio},\texttt{Sales})$ captures is already being explained by \texttt{TV}. This is exactly why \textbf{multiple regression} (Section 09.07) --- rather than a collection of isolated bivariate correlations --- is the correct tool for quantifying each predictor's effect \emph{controlling for the others}.
  \end{alertblock}
\end{frame}

% Slide 14: Closing
\begin{frame}{Section Closing: Correlation}
  \begin{reflexion}
    Correlation quantifies the strength and direction of a linear association, and is the conceptual motivation for the whole chapter: if two variables are unrelated, no regression model can usefully predict one from the other. But correlation describes association, not causal mechanism --- a distinction every data scientist must keep in mind.
  \end{reflexion}

  \vspace{0.2cm}
  \pause
  \begin{block}{Key Takeaways}
    \begin{itemize}\itemsep=0.08cm
      \item \textbf{Range:} $r_{XY}\in[-1,1]$, with sign indicating direction and magnitude indicating strength.
      \item \textbf{Correlation $\neq$ causation:} a confounding third variable can generate spurious association.
      \item \textbf{Bivariate correlation $\neq$ controlled effect:} predictors correlated with each other require multiple regression, not isolated correlations.
    \end{itemize}
  \end{block}
\end{frame}

% Slide 15: Modular Perspective
\begin{frame}{Modular Perspective --- The Next Challenge}
  \scriptsize
  \begin{retoclase}
    Correlation measures association but does not let us predict a specific numeric value of $Y$ given a value of $X$. Section 09.02 introduces \textbf{linear regression}: the model $Y_e=\hat\beta_0+\hat\beta_1X$ that turns the association measured by $r_{XY}$ into a concrete predictive tool.
  \end{retoclase}

  \vspace{0.08cm}
  \pause
  \begin{alertblock}{Recommended Preparation}
    \begin{itemize}\itemsep=0.06cm
      \item Reflect on what additional information (beyond $r_{XY}$) you would need to predict a specific \texttt{Sales} value given a level of \texttt{TV} spending.
      \item Review the equation of a line $y=ax+b$: linear regression will estimate $a$ and $b$ from the data.
    \end{itemize}
  \end{alertblock}
\end{frame}

\end{document}
